% !TEX root = ../main.tex
\begin{table}[htbp]
  \centering
  \caption[Efecto del tamaño máximo con la estrategia elegida]{Efecto
      del tamaño máximo del fragmento, manteniendo fija la estrategia estructural con el
      objetivo igualado al máximo.}
  \label{tab:tamano_fragmento}
  \small
  \begin{tabular}{@{}rrrrrrr@{}}
    \toprule
    \textbf{Máx.} & \textbf{Frag.} & \textbf{Mediana} & \textbf{RU@1}  & \textbf{RU@3}  & \textbf{MRR}   & \textbf{Trunc.} \\
    \midrule
    600           & 2.695          & 569              & 0,930          & 0,980          & 0,960          & 0               \\
    \textbf{900}  & \textbf{1.334} & \textbf{838}     & \textbf{0,930} & \textbf{0,985} & \textbf{0,970} & \textbf{0}      \\
    1.200         & 947            & 1.104            & 0,850          & 0,960          & 0,913          & 0               \\
    1.500         & 753            & 1.344            & 0,830          & 0,945          & 0,898          & 4               \\
    1.800         & 632            & 1.497            & 0,770          & 0,945          & 0,868          & 11              \\
    \bottomrule
  \end{tabular}
  \begin{minipage}{0.95\textwidth}
    \vspace{4pt}
    \tiny
    \textbf{Notas:} \textbf{Máx.}: tamaño máximo del fragmento, en caracteres, que es una
    restricción dura; \textbf{Frag.}: número total de fragmentos que produce esa
    configuración sobre la colección completa; \textbf{Mediana}: tamaño mediano del
    fragmento resultante, en caracteres; \textbf{RU@K} (\textit{Recall por Unidad}):
    proporción de preguntas cuya asignatura correcta aparece entre los $K$ primeros
    resultados; \textbf{MRR} (\textit{Mean Reciprocal Rank}): media del inverso de la
    posición del primer resultado correcto; \textbf{Trunc.}: fragmentos que superan la
    ventana del modelo de incrustaciones, contados con su propio analizador léxico, y que el
    modelo recorta en silencio, sin avisar ni fallar. En negrita, la configuración adoptada.
  \end{minipage}
\end{table}
